% BHU PYQ -- Chemistry: Physical Chemistry-I (2025-26 NEP)
% Paper Code: CHEMJ33 (Sem III Mid-Semester Examination, 2025-26)
\documentclass[11pt,a4paper]{article}
\usepackage[a4paper,top=2.5cm,bottom=2.5cm,left=2.8cm,right=2.8cm,headheight=14pt]{geometry}
\usepackage{amsmath,amssymb}
\usepackage[T1]{fontenc}
\usepackage[utf8]{inputenc}
\usepackage{lmodern,microtype,fancyhdr,enumitem,hyperref,lastpage,parskip}

\hypersetup{
    colorlinks=true,
    urlcolor=black,
    linkcolor=black,
    pdftitle={CHEMJ33: Physical Chemistry-I -- BHU Sem III Mid-Term 2025-26 (NEP)}
}

\pagestyle{fancy}\fancyhf{}
\renewcommand{\headrulewidth}{0.4pt}\renewcommand{\footrulewidth}{0.4pt}
\fancyhead[L]{\small \textit{Banaras Hindu University}}
\fancyhead[R]{\small \textit{CHEMJ33: Physical Chemistry-I (Mid-Term 2025-26)}}
\fancyfoot[C]{\small Page \thepage\ of \pageref{LastPage}}

\newcommand{\pts}[1]{\hfill \textit{[#1]}}

\begin{document}

\begin{center}
  {\large\bfseries BANARAS HINDU UNIVERSITY}\\[2pt]
  {\normalsize B.Sc. (Hons.) Chemistry Semester-III Mid-Semester Examination 2025-26 (NEP)}\\[6pt]
  \rule{0.8\linewidth}{0.5pt}\\[6pt]
  {\large\bfseries CHEMISTRY}\\[4pt]
  {\large\bfseries Paper Code: CHEMJ33 : Physical Chemistry-I}\\[4pt]
  {\normalsize Date: 04.11.2025 \hfill Time: 1 Hour \hfill Maximum Marks: 20}
\end{center}

\rule{\linewidth}{0.4pt}
\smallskip

\noindent \textit{Instructions: Answer all the questions. Figures in the right-hand margin indicate marks.}

\bigskip

\begin{enumerate}[label=\textbf{\arabic*.},leftmargin=*,itemsep=12pt]

  \item Define molar conductance and equivalent conductance of an electrolyte. Why does the conductivity of a solution decrease with dilution? \pts{2}

  \item Calculate the degree of dissociation ($\alpha$) of acetic acid if its molar conductivity ($\Lambda_m$) is $39.05\text{ S cm}^2\text{ mol}^{-1}$. Given: $\lambda^\circ(\text{H}^+) = 349.6\text{ S cm}^2\text{ mol}^{-1}$ and $\lambda^\circ(\text{CH}_3\text{COO}^-) = 40.9\text{ S cm}^2\text{ mol}^{-1}$. \pts{2}

  \item Define transference (transport) number of ions. Write the relation between transference number and ionic mobility. \pts{2}

  \item Plot and explain a graph for the conductometric titration of weak acid by strong base. \pts{2}

  \item A moving boundary experiment is conducted to measure the transference number of the $\text{Li}^+$ ion in a $0.01\text{ mol L}^{-1}$ $\text{LiCl}$ solution. In a tube having a cross-sectional area of $0.125\text{ cm}^2$, the boundary moves $7.3\text{ cm}$ in $1490\text{ sec}$ using a current of $1.8 \times 10^{-3}\text{ A}$. Calculate the transference numbers of the cation and anion. \pts{2}

  \item Define cell potential. How can you determine the standard electrode potential of an electrode? \pts{2}

  \item What is liquid junction potential? How does a salt bridge minimise the liquid junction potential of a cell? \pts{2}

  \item For the concentration cell:
  \[
    \text{Pt} \mid \text{H}_2(\text{g}, 1\text{ atm}) \mid \text{HCl soln. } (a_1) : \text{HCl soln. } (a_2) \mid \text{H}_2(\text{g}, 1\text{ atm}) \mid \text{Pt}
  \]
  write the various processes at the two electrodes and at the liquid junction. \pts{3}

  \item A $50\text{ mL}$ $0.05\text{ M}$ solution of $\text{Fe(II)}$ is titrated with $0.05\text{ M}$ solution of $\text{Ce(IV)}$ in the presence of dilute $\text{H}_2\text{SO}_4$ at $25^\circ\text{C}$ by coupling with saturated calomel electrode. Draw the potentiometric curve ($E_{\text{cell}}$ vs. volume of $\text{Ce(IV)}$) and calculate the equivalence point potential. [Given, $E^\circ_{\text{Fe}^{3+}/\text{Fe}^{2+}} = 0.75\text{ V}$ and $E^\circ_{\text{Ce}^{4+}/\text{Ce}^{3+}} = 1.45\text{ V}$ at $25^\circ\text{C}$] \pts{3}

\end{enumerate}

\bigskip
\begin{center}
  \textbf{***}
\end{center}

\end{document}
